% Options for packages loaded elsewhere
% Options for packages loaded elsewhere
\PassOptionsToPackage{unicode}{hyperref}
\PassOptionsToPackage{hyphens}{url}
%
\documentclass[
  english,
  russian,
  12pt,
  a4paper,
  DIV=11,
  numbers=noendperiod]{scrreprt}
\usepackage{xcolor}
\usepackage{amsmath,amssymb}
\setcounter{secnumdepth}{5}
\usepackage{iftex}
\ifPDFTeX
  \usepackage[T1]{fontenc}
  \usepackage[utf8]{inputenc}
  \usepackage{textcomp} % provide euro and other symbols
\else % if luatex or xetex
  \usepackage{unicode-math} % this also loads fontspec
  \defaultfontfeatures{Scale=MatchLowercase}
  \defaultfontfeatures[\rmfamily]{Ligatures=TeX,Scale=1}
\fi
\usepackage{lmodern}
\ifPDFTeX\else
  % xetex/luatex font selection
\fi
% Use upquote if available, for straight quotes in verbatim environments
\IfFileExists{upquote.sty}{\usepackage{upquote}}{}
\IfFileExists{microtype.sty}{% use microtype if available
  \usepackage[]{microtype}
  \UseMicrotypeSet[protrusion]{basicmath} % disable protrusion for tt fonts
}{}
\usepackage{setspace}
% Make \paragraph and \subparagraph free-standing
\makeatletter
\ifx\paragraph\undefined\else
  \let\oldparagraph\paragraph
  \renewcommand{\paragraph}{
    \@ifstar
      \xxxParagraphStar
      \xxxParagraphNoStar
  }
  \newcommand{\xxxParagraphStar}[1]{\oldparagraph*{#1}\mbox{}}
  \newcommand{\xxxParagraphNoStar}[1]{\oldparagraph{#1}\mbox{}}
\fi
\ifx\subparagraph\undefined\else
  \let\oldsubparagraph\subparagraph
  \renewcommand{\subparagraph}{
    \@ifstar
      \xxxSubParagraphStar
      \xxxSubParagraphNoStar
  }
  \newcommand{\xxxSubParagraphStar}[1]{\oldsubparagraph*{#1}\mbox{}}
  \newcommand{\xxxSubParagraphNoStar}[1]{\oldsubparagraph{#1}\mbox{}}
\fi
\makeatother


\usepackage{longtable,booktabs,array}
\usepackage{calc} % for calculating minipage widths
% Correct order of tables after \paragraph or \subparagraph
\usepackage{etoolbox}
\makeatletter
\patchcmd\longtable{\par}{\if@noskipsec\mbox{}\fi\par}{}{}
\makeatother
% Allow footnotes in longtable head/foot
\IfFileExists{footnotehyper.sty}{\usepackage{footnotehyper}}{\usepackage{footnote}}
\makesavenoteenv{longtable}
\usepackage{graphicx}
\makeatletter
\newsavebox\pandoc@box
\newcommand*\pandocbounded[1]{% scales image to fit in text height/width
  \sbox\pandoc@box{#1}%
  \Gscale@div\@tempa{\textheight}{\dimexpr\ht\pandoc@box+\dp\pandoc@box\relax}%
  \Gscale@div\@tempb{\linewidth}{\wd\pandoc@box}%
  \ifdim\@tempb\p@<\@tempa\p@\let\@tempa\@tempb\fi% select the smaller of both
  \ifdim\@tempa\p@<\p@\scalebox{\@tempa}{\usebox\pandoc@box}%
  \else\usebox{\pandoc@box}%
  \fi%
}
% Set default figure placement to htbp
\def\fps@figure{htbp}
\makeatother



\ifLuaTeX
\usepackage[bidi=basic,provide=*]{babel}
\else
\usepackage[bidi=default,provide=*]{babel}
\fi
% get rid of language-specific shorthands (see #6817):
\let\LanguageShortHands\languageshorthands
\def\languageshorthands#1{}


\setlength{\emergencystretch}{3em} % prevent overfull lines

\providecommand{\tightlist}{%
  \setlength{\itemsep}{0pt}\setlength{\parskip}{0pt}}



 
\usepackage[style=gost-numeric,backend=biber,langhook=extras,autolang=other*]{biblatex}
\addbibresource{bib/cite.bib}

\usepackage[]{csquotes}

\usepackage{indentfirst}
\usepackage{float}
\floatplacement{figure}{H}
\IfFileExists{plex-otf.sty}{
  %% Full TeXlive
  \usepackage[
    % math,
    RM={Scale=0.94},SS={Scale=0.94},SScon={Scale=0.94},TT={Scale=MatchLowercase,FakeStretch=0.9},
    DefaultFeatures={Ligatures=Common}
  ]{plex-otf}
  % \usepackage[Scale=MatchUppercase]{juliamono}
}{
  %% TinyTeX
  \usepackage{libertine}
}
\KOMAoption{captions}{tableheading}
\makeatletter
\@ifpackageloaded{caption}{}{\usepackage{caption}}
\AtBeginDocument{%
\ifdefined\contentsname
  \renewcommand*\contentsname{Содержание}
\else
  \newcommand\contentsname{Содержание}
\fi
\ifdefined\listfigurename
  \renewcommand*\listfigurename{Список иллюстраций}
\else
  \newcommand\listfigurename{Список иллюстраций}
\fi
\ifdefined\listtablename
  \renewcommand*\listtablename{Список таблиц}
\else
  \newcommand\listtablename{Список таблиц}
\fi
\ifdefined\figurename
  \renewcommand*\figurename{Рисунок}
\else
  \newcommand\figurename{Рисунок}
\fi
\ifdefined\tablename
  \renewcommand*\tablename{Таблица}
\else
  \newcommand\tablename{Таблица}
\fi
}
\@ifpackageloaded{float}{}{\usepackage{float}}
\floatstyle{ruled}
\@ifundefined{c@chapter}{\newfloat{codelisting}{h}{lop}}{\newfloat{codelisting}{h}{lop}[chapter]}
\floatname{codelisting}{Список}
\newcommand*\listoflistings{\listof{codelisting}{Листинги}}
\makeatother
\makeatletter
\makeatother
\makeatletter
\@ifpackageloaded{caption}{}{\usepackage{caption}}
\@ifpackageloaded{subcaption}{}{\usepackage{subcaption}}
\makeatother
\usepackage{bookmark}
\IfFileExists{xurl.sty}{\usepackage{xurl}}{} % add URL line breaks if available
\urlstyle{same}
\hypersetup{
  pdftitle={Отчёт по лабораторной работе №2},
  pdfauthor={Полина Вячеславовна Белакова},
  pdflang={ru-RU},
  hidelinks,
  pdfcreator={LaTeX via pandoc}}


\title{Отчёт по лабораторной работе №2}
\author{Полина Вячеславовна Белакова}
\date{}
\begin{document}
\maketitle

\renewcommand*\contentsname{Содержание}
{
\setcounter{tocdepth}{1}
\tableofcontents
}
\listoffigures
\listoftables

\setstretch{1.5}
\chapter{Цель
работы}\label{ux446ux435ux43bux44c-ux440ux430ux431ux43eux442ux44b}

Изучить идеологию и применение средств контроля версий и приобрести
навыки работы с git.

\chapter{Задание}\label{ux437ux430ux434ux430ux43dux438ux435}

\begin{itemize}
\tightlist
\item
  Изучить идеологию и применение средств контроля версий.
\item
  Освоить умения по работе с git.
\end{itemize}

\chapter{Теоретическое
введение}\label{ux442ux435ux43eux440ux435ux442ux438ux447ux435ux441ux43aux43eux435-ux432ux432ux435ux434ux435ux43dux438ux435}

Системы контроля версий (Version Control System, VCS) применяются при
работе нескольких человек над одним проектом. Обычно основное дерево
проекта хранится в локальном или удалённом репозитории, к которому
настроен доступ для участников проекта. При внесении изменений в
содержание проекта система контроля версий позволяет их фиксировать,
совмещать изменения, произведённые разными участниками проекта,
производить откат к любой более ранней версии проекта, если это
требуется. Среди классических VCS наиболее известны CVS, Subversion, а
среди распределённых --- Git, Bazaar, Mercurial. Принципы их работы
схожи, отличаются они в основном синтаксисом используемых в работе
команд. Система контроля версий Git представляет собой набор программ
командной строки. Доступ к ним можно получить из терминала посредством
ввода команды git с различными опциями. Благодаря тому, что Git является
распределённой системой контроля версий, резервную копию локального
хранилища можно сделать простым копированием или архивацией.

\chapter{Выполнение лабораторной
работы}\label{ux432ux44bux43fux43eux43bux43dux435ux43dux438ux435-ux43bux430ux431ux43eux440ux430ux442ux43eux440ux43dux43eux439-ux440ux430ux431ux43eux442ux44b}

\section{Установка программного
обеспечения}\label{ux443ux441ux442ux430ux43dux43eux432ux43aux430-ux43fux440ux43eux433ux440ux430ux43cux43cux43dux43eux433ux43e-ux43eux431ux435ux441ux43fux435ux447ux435ux43dux438ux44f}

Установила git(рис.~\ref{fig-001}).

\begin{figure}

\centering{

\includegraphics[width=0.7\linewidth,height=\textheight,keepaspectratio]{image/install_git1.png}

}

\caption{\label{fig-001}Установка git}

\end{figure}%

Установила gh (рис.~\ref{fig-002}).

\begin{figure}

\centering{

\includegraphics[width=0.7\linewidth,height=\textheight,keepaspectratio]{image/install_gh2.png}

}

\caption{\label{fig-002}Установка gh}

\end{figure}%

\section{Базовая настройка
git}\label{ux431ux430ux437ux43eux432ux430ux44f-ux43dux430ux441ux442ux440ux43eux439ux43aux430-git}

Зададаю имя и email владельца репозитория (рис.~\ref{fig-003}).

\begin{figure}

\centering{

\includegraphics[width=0.7\linewidth,height=\textheight,keepaspectratio]{image/git_config_user_email3.png}

}

\caption{\label{fig-003}Настройка имени и почты владельца репозитория}

\end{figure}%

Настраиваю utf-8 в выводе сообщений, верификацию и подписание коммитов
git, зададаю имя начальной ветки, настраиваю параметры autocrlf,
safecrlf Установила gh (рис.~\ref{fig-004}).

\begin{figure}

\centering{

\includegraphics[width=0.7\linewidth,height=\textheight,keepaspectratio]{image/configi5.png}

}

\caption{\label{fig-004}Настройка git}

\end{figure}%

Создаю ключи ssh (рис.~\ref{fig-005}).

\begin{figure}

\centering{

\includegraphics[width=0.7\linewidth,height=\textheight,keepaspectratio]{image/ssh6.png}

}

\caption{\label{fig-005}Создание ключа ssh}

\end{figure}%

Создаю ключи pgp(рис.~\ref{fig-006}).

\begin{figure}

\centering{

\includegraphics[width=0.7\linewidth,height=\textheight,keepaspectratio]{image/pgp7.png}

}

\caption{\label{fig-006}Создание ключа pgp}

\end{figure}%

\section{Настройка
github}\label{ux43dux430ux441ux442ux440ux43eux439ux43aux430-github}

Добавляю PGP ключ в GitHub(рис.~\ref{fig-007}).

\begin{figure}

\centering{

\includegraphics[width=0.7\linewidth,height=\textheight,keepaspectratio]{image/gpg_list8.png}

}

\caption{\label{fig-007}Добавление ключа pgp}

\end{figure}%

Копирую сгенерированный PGP ключ в буфер обмена (рис.~\ref{fig-008}).

\begin{figure}

\centering{

\includegraphics[width=0.7\linewidth,height=\textheight,keepaspectratio]{image/gpg_armore.png}

}

\caption{\label{fig-008}Копирование ключа pgp}

\end{figure}%

Перехожу в настройки GitHub, нажимаю на кнопку New GPG key и вставьте
полученный ключ в поле ввода(\textbf{?@fig-009}).

\includegraphics[width=0.7\linewidth,height=\textheight,keepaspectratio]{image/add.png}.

\section{Настройка автоматических подписей коммитов
git}\label{ux43dux430ux441ux442ux440ux43eux439ux43aux430-ux430ux432ux442ux43eux43cux430ux442ux438ux447ux435ux441ux43aux438ux445-ux43fux43eux434ux43fux438ux441ux435ux439-ux43aux43eux43cux43cux438ux442ux43eux432-git}

Используя введёный email, указываю примениние Git при подписи коммитов
(рис.~\ref{fig-010}).

\begin{figure}

\centering{

\includegraphics[width=0.7\linewidth,height=\textheight,keepaspectratio]{image/git_configi.png}

}

\caption{\label{fig-010}Настраиваю автоматическую подпись коммитов}

\end{figure}%

\section{Настройка
gh}\label{ux43dux430ux441ux442ux440ux43eux439ux43aux430-gh}

Авторизоввываюсь с помощью команды gh auth login (рис.~\ref{fig-011}).

\begin{figure}

\centering{

\includegraphics[width=0.7\linewidth,height=\textheight,keepaspectratio]{image/auth_login.png}

}

\caption{\label{fig-011}Авторизация}

\end{figure}%

\section{Шаблон для рабочего
пространства}\label{ux448ux430ux431ux43bux43eux43d-ux434ux43bux44f-ux440ux430ux431ux43eux447ux435ux433ux43e-ux43fux440ux43eux441ux442ux440ux430ux43dux441ux442ux432ux430}

Создаю дерикторию (рис.~\ref{fig-012})и репозиторий курса на основе
шаблона(рис.~\ref{fig-013}).

\begin{figure}

\centering{


\includegraphics[width=0.7\linewidth,height=\textheight,keepaspectratio]{image/mkdir.png}

}

\caption{\label{fig-012}Создание директории}

\end{figure}%

\begin{figure}

\centering{

\includegraphics[width=0.7\linewidth,height=\textheight,keepaspectratio]{image/createclone.png}

}

\caption{\label{fig-013}Создание репозиторя и его клонирование}

\end{figure}%

\section{Настройка каталога
курса}\label{ux43dux430ux441ux442ux440ux43eux439ux43aux430-ux43aux430ux442ux430ux43bux43eux433ux430-ux43aux443ux440ux441ux430}

Перехожу в каталог курса, удаляю лишние файлы (рис.~\ref{fig-014}),
создаю необходимые каталоги (рис.~\ref{fig-015})

\begin{figure}

\centering{

\includegraphics[width=0.7\linewidth,height=\textheight,keepaspectratio]{image/cdlsrmls.png}

}

\caption{\label{fig-014}Удаление лишних файлов}

\end{figure}%

\begin{figure}

\centering{

\includegraphics[width=0.7\linewidth,height=\textheight,keepaspectratio]{image/echomake.png}

}

\caption{\label{fig-015}Создание необходимых каталогов}

\end{figure}%

Отправляю файлы на сервер (рис.~\ref{fig-016}), (рис.~\ref{fig-017}).

\begin{figure}

\centering{

\includegraphics[width=0.7\linewidth,height=\textheight,keepaspectratio]{image/addcommit.png}

}

\caption{\label{fig-016}Выполнение git add и git commit}

\end{figure}%

\begin{figure}

\centering{

\includegraphics[width=0.7\linewidth,height=\textheight,keepaspectratio]{image/git_push.png}

}

\caption{\label{fig-017}Выполнение git push}

\end{figure}%

\chapter{Выводы}\label{ux432ux44bux432ux43eux434ux44b}

Была изучена идеология и применение средств контроля версий.Приобретены
навыки работы с git.

\chapter{Контрольные
вопросы}\label{ux43aux43eux43dux442ux440ux43eux43bux44cux43dux44bux435-ux432ux43eux43fux440ux43eux441ux44b}

\begin{enumerate}
\def\labelenumi{\arabic{enumi}.}
\item
  Что такое системы контроля версий (VCS) и для решения каких задач они
  предназначаются? Система контроля версий (VCS) --- это программное
  обеспечение для отслеживания изменений в файлах и координации работы
  над ними. VCS хранит историю изменений, позволяет возвращаться к
  предыдущим версиям и определять авторство каждой правки. Главные
  задачи: обеспечение совместной работы нескольких человек над проектом,
  защита от потери данных и возможность экспериментировать без риска
  сломать основной код.
\item
  Объясните понятия : хранилище, commit, история, рабочая копия.
  Хранилище (repository) --- это место, где VCS хранит все файлы проекта
  и всю историю их изменений. Commit --- операция сохранения текущего
  состояния файлов в хранилище, \enquote{снимок} проекта в определённый
  момент времени. История --- цепочка последовательных commit'ов,
  позволяющая проследить эволюцию проекта. Рабочая копия --- текущая
  версия файлов на вашем компьютере, с которой вы непосредственно
  работаете.
\item
  Чем отличаются централизованные и децентрализованные VCS? В
  централизованных VCS (например, SVN) существует одно главное хранилище
  на сервере, а разработчики получают только рабочие копии; при
  недоступности сервера работа с историей невозможна. В
  децентрализованных VCS (например, Git) каждый разработчик имеет полную
  копию хранилища с историей на своём компьютере, что позволяет работать
  автономно. Любая локальная копия в DVCS является полноценным резервным
  хранилищем проекта.
\item
  Опишите действия с VCS при единоличной работе. При единоличной работе
  сначала инициализируется пустое хранилище в папке проекта (git init).
  Затем создаются и изменяются файлы, после чего они добавляются в
  индекс (git add) и фиксируются (git commit) с пояснительным
  сообщением. В любой момент можно посмотреть историю изменений (git
  log) или вернуться к более ранней версии проекта.
\item
  Опишите порядок работы с общим хранилищем VCS. Работа начинается с
  клонирования удалённого репозитория (git clone). Перед началом работы
  и перед отправкой изменений необходимо загрузить актуальную версию от
  коллег (git pull). После внесения изменений они фиксируются локально
  (git commit), а затем отправляются в общее хранилище (git push). При
  возникновении конфликтов их нужно разрешить вручную и сделать commit
  слияния.
\item
  Каковы основные задачи git? Git обеспечивает полную историю изменений
  проекта с возможностью навигации по ней. Он поддерживает нелинейную
  разработку через мощную систему ветвления и гарантирует целостность
  данных с помощью контрольных сумм. Git позволяет каждому разработчику
  иметь полную копию репозитория и работать автономно, а также
  предоставляет гибкие инструменты для управления версиями (теги,
  ветки).
\item
  Назовите и дайте краткую характеристику командам git. git init создаёт
  новый локальный репозиторий, а git clone копирует удалённый. git add
  добавляет изменения в индекс, git commit фиксирует их в репозитории.
  git push отправляет изменения на удалённый сервер, а git pull
  загружает их оттуда. git branch и git checkout управляют ветками, git
  merge сливает ветки, а git status и git log показывают состояние и
  историю.
\item
  Приведите примеры использования с локальным и удалённым репозиториями.
  Локально git init → создаём файлы → git add → git commit -m
  'сообщение'' → git log для просмотра истории. С удалённым: git clone →
  git pull для обновления → работаем → git commit → git push для
  отправки изменений. Для новой задачи создаём ветку git checkout -b
  feature, а после завершения отправляем её git push origin feature.
\item
  Что такое ветви (branches) и зачем они нужны? Ветка --- это
  перемещаемый указатель на один из commit'ов, позволяющий вести
  разработку изолированно от основной линии. Ветки нужны для
  параллельной разработки новых функций без риска сломать стабильный
  код. Они также используются для срочного исправления ошибок,
  проведения экспериментов и организации работы нескольких разработчиков
  над одним проектом.
\item
  Как и зачем игнорировать файлы при commit? Для игнорирования файлов
  используется файл .gitignore, в котором перечисляются шаблоны имён
  файлов и папок, не подлежащих версионированию. Игнорируют временные
  файлы редакторов, скомпилированные бинарные файлы, конфиденциальные
  данные и папки с зависимостями. Это позволяет держать репозиторий
  чистым и лёгким, содержащим только исходный код и документацию.
\end{enumerate}

\chapter*{Список
литературы}\label{ux441ux43fux438ux441ux43eux43a-ux43bux438ux442ux435ux440ux430ux442ux443ux440ux44b}
\addcontentsline{toc}{chapter}{Список литературы}

\printbibliography[heading=none]





\end{document}
